% Źródła: Eves s. 287; Wikipedia: János Bolyai, Farkas Bolyai, Carl Friedrich Gauss, Non-Euclidean geometry.
% https://en.wikipedia.org/wiki/J%C3%A1nos_Bolyai
% https://en.wikipedia.org/wiki/Farkas_Bolyai
% https://www.deltami.edu.pl/2012/06/w-rozumowaniach-byl-blad/
% https://www.deltami.edu.pl/2018/08/geometria-bolyaia-lobaczewskiego-co-to-jest-i-jak-ja-poznawac/

\section{Gauss oraz Farkas i János Bolyai} % lang-pl
\section{Gauss e Farkas e János Bolyai} % lang-it
\label{sekcja_gauss_bolyai}

Dwaj przyjaciele ze studiów i syn jednego z nich odkryją tę samą geometrię, a jeden z nich nic o tym nie opublikuje. % lang-pl
W tej sekcji zobaczymy, kto komu napisze i dlaczego pierwszeństwo będzie sprawą tak drażliwą. % lang-pl
Due amici di studi e il figlio di uno dei due scopriranno la stessa geometria, e uno di loro non ne pubblicherà nulla. % lang-it
In questa sezione vedremo chi scriverà a chi e perché la priorità sarà una questione così delicata. % lang-it

Carl Friedrich Gauss (1777--1855) zapozna się w Getyndze, około 1796--1799 roku, z Farkasem Bolyaiem, i ci dwaj zostaną bliskimi przyjaciółmi. % lang-pl
Farkas Bolyai (1775--1856), urodzony w Bólyi, a zmarły w Marosvásárhely, przez całe życie będzie zajmował się podstawami geometrii i aksjomatem równoległych. % lang-pl
Carl Friedrich Gauss (1777--1855) conoscerà a Gottinga, intorno al 1796--1799, Farkas Bolyai, e i due diventeranno amici intimi. % lang-it
Farkas Bolyai (1775--1856), nato a Bólya e morto a Marosvásárhely, si occuperà per tutta la vita dei fondamenti della geometria e dell'assioma delle parallele. % lang-it
\index[persons]{Gauss, Carl Friedrich}%
\index[persons]{Bolyai, Farkas}%

Jego syn János (1802--1860), urodzony w Kolozsvárze, mimo ostrzeżeń ojca weźmie się za ten sam problem. % lang-pl
Farkas napisze mu, że przeszedł przez tę „bezdenną noc'', która zgasiła w jego życiu wszelkie światło i radość. % lang-pl
János nie usłucha, a w 1823 roku napisze do ojca z ekscytacją, że odkrył rzeczy tak cudowne, iż sam się zdumiał, i że „z niczego stworzył dziwny nowy świat''. % lang-pl
Suo figlio János (1802--1860), nato a Kolozsvár, nonostante gli avvertimenti del padre si dedicherà allo stesso problema. % lang-it
Farkas gli scriverà di aver attraversato quella «notte senza fondo» che aveva spento nella sua vita ogni luce e ogni gioia. % lang-it
János non gli darà retta e nel 1823 scriverà al padre, entusiasta, di aver scoperto cose tanto meravigliose da restarne stupefatto, e di aver «creato dal nulla uno strano nuovo universo». % lang-it
\index[persons]{Bolyai, János}%

W 1832 roku wyniki Jánosa ukażą się jako dodatek do pracy ojca \emph{Tentamen}: Farkas, początkowo niechętny, dojdzie do entuzjazmu i sam przekona syna do publikacji. % lang-pl
Gauss odpowie z uznaniem, nazywając Jánosa geniuszem pierwszej rangi, ale doda, że sam doszedł do podobnych wyników wiele lat wcześniej -- co, jak się mówi, wywoła napięcie między ojcem i synem. % lang-pl
Później János dowie się, że już w latach 1829--1830 podobne wyniki opublikował Łobaczewski (zobacz sekcję \ref{sekcja_lobaczewski}). % lang-pl
Nel 1832 i risultati di János usciranno come appendice all'opera del padre, il \emph{Tentamen}: Farkas, dapprima riluttante, si entusiasmerà e convincerà lui stesso il figlio a pubblicare. % lang-it
Gauss risponderà con ammirazione, definendo János un genio di prim'ordine, ma aggiungerà di essere giunto a risultati simili molti anni prima -- cosa che, si dice, creerà tensione tra padre e figlio. % lang-it
In seguito János verrà a sapere che già nel 1829--1830 risultati analoghi erano stati pubblicati da Lobachevskij (si veda la sezione \ref{sekcja_lobaczewski}). % lang-it

W 1829 roku János Bolyai, Carl Friedrich Gauss oraz (niezależnie od nich) Nikołaj Łobaczewski stworzą geometrię hiperboliczną, czyli pierwszą w pełni rozwiniętą geometrię nieeuklidesową, w której piąty postulat Euklidesa nie obowiązuje. % lang-pl
Ściślej: Łobaczewski opublikuje swoje wyniki w latach 1829--1830, Bolyai w 1832 roku, a Gauss -- choć według Evesa pierwszy dojdzie do głębokich wniosków -- nie opublikuje niczego, więc honor odkrycia trzeba będzie z nimi podzielić \cite[s. 287]{eves1_1972}. % lang-pl
Nel 1829 János Bolyai, Carl Friedrich Gauss e, indipendentemente da loro, Nikolaj Lobachevskij costruiranno la geometria iperbolica, la prima geometria non euclidea pienamente sviluppata, nella quale il quinto postulato di Euclide non vale. % lang-it
Più esattamente: Lobachevskij pubblicherà i suoi risultati nel 1829--1830, Bolyai nel 1832, e Gauss -- pur essendo, secondo Eves, il primo a giungere a conclusioni profonde -- non pubblicherà nulla, sicché l'onore della scoperta dovrà essere diviso con loro \cite[p. 287]{eves1_1972}. % lang-it

Gauss, według Evesa, podejdzie do problemu tak samo jak pozostali dwaj: przez trzy możliwości dla aksjomatu Playfaira -- przez punkt przechodzi więcej niż jedna, dokładnie jedna albo żadna prosta równoległa do danej -- a trzecią łatwo wykluczy przy założeniu nieskończoności prostej \cite[s. 287]{eves1_1972}. % lang-pl
On pierwszy nada nazwę \emph{geometria nieeuklidesowa}, początkowo mając na myśli własną geometrię hiperboliczną; w 1824 roku napisze też list do Taurinusa (zobacz sekcję \ref{sekcja_schweikart_taurinus}). % lang-pl
Gauss, secondo Eves, affronterà il problema come gli altri due: attraverso le tre possibilità dell'assioma di Playfair -- per un punto passa più di una, esattamente una o nessuna retta parallela a una data -- e la terza la escluderà facilmente supponendo infinita la retta \cite[p. 287]{eves1_1972}. % lang-it
Sarà lui il primo a chiamare \emph{geometria non euclidea} la propria geometria iperbolica; nel 1824 scriverà anche una lettera a Taurinus (si veda la sezione \ref{sekcja_schweikart_taurinus}). % lang-it

János Bolyai zbuduje geometrię, która obejmuje jednocześnie euklidesową i hiperboliczną, zależnie od parametru $k$; nazwie ją \emph{geometrią absolutną}. % lang-pl
Zauważy też, że samym rozumowaniem matematycznym nie da się rozstrzygnąć, czy geometria świata fizycznego jest euklidesowa, czy nie (zobacz sekcję \ref{sekcja_fizyka}). % lang-pl
Eves nazywa geometrią absolutną wszystko, co da się udowodnić bez piątego postulatu -- czyli naszą geometrię neutralną; do niej należy pierwsze dwadzieścia osiem twierdzeń pierwszej księgi Euklidesa \cite[s. 288]{eves1_1972}. % lang-pl
János Bolyai costruirà una geometria che comprende insieme quella euclidea e quella iperbolica, a seconda di un parametro $k$; la chiamerà \emph{geometria assoluta}. % lang-it
Noterà anche che con il solo ragionamento matematico non si può decidere se la geometria del mondo fisico sia euclidea o meno (si veda la sezione \ref{sekcja_fizyka}). % lang-it
Eves chiama geometria assoluta tutto ciò che si dimostra senza il quinto postulato -- cioè la nostra geometria neutrale; ne fanno parte le prime ventotto proposizioni del primo libro di Euclide \cite[p. 288]{eves1_1972}. % lang-it
\index{geometria!absolutna} % lang-pl
\index{geometria!assoluta} % lang-it

% Uwaga: nie opisuję listów Gaussa (do Taurinusa 1824, do Farkasa 1832) dokładnie -- potrzebne źródło z datami.
